\chapter{Methodology and Design}
\label{ch:methodology}

FoodLink was built iteratively~\cite{foodlinkProposal,sommerville2016}, as small
numbered tasks (the latest is Task~46), each with its own regression tests, and
every non-trivial choice was recorded in a decision log cited below by number
(D-01 to D-61)~\cite{foodlinkArchitecture}. Three audits of the running system
(2, 5 and 10 September 2026) ranked findings from P0 to P3; all five P1 findings
are fixed.

\section{System Architecture}
\label{sec:architecture}

\subsection{High-Level Design}
\label{sec:hld}

FoodLink has three tiers: a React single-page application, a stateless REST API
built with FastAPI, and a relational database (SQLite by default) reached only
through the SQLAlchemy ORM. The API is the sole authority for every rule, makes
no outbound requests (no email, SMS, maps or object storage) and runs no
background workers. Five routers (\code{auth}, \code{donations},
\code{organisations}, \code{admin}, \code{metrics}) call shared modules with no
HTTP logic (\code{matching}, \code{serialize}, \code{security},
\code{ratelimit}); there is deliberately no service layer (D-07).

\subsection{Technical Stack}
\label{sec:stack}

\begin{table}[!htbp]
\centering
\small
\begin{tabular}{|L{2.6cm}|L{10.7cm}|}
  \hline
  \textbf{Layer} & \textbf{Technologies} \\
  \hline
  Backend & Python 3.13 (3.10 or later required); FastAPI $\geq$\,0.115~\cite{fastapiDocs}, uvicorn, Pydantic~2, PyJWT, bcrypt \\
  \hline
  Data & SQLAlchemy 2.0~\cite{sqlalchemyDocs}; Alembic $\geq$\,1.13; SQLite by default~\cite{sqliteWhenToUse} \\
  \hline
  Frontend & React 18.3~\cite{reactDocs}, TypeScript 5.5, Vite 5.4, React Router 6, Tailwind CSS 3.4 \\
  \hline
  Testing and CI & pytest with FastAPI's test client; Vitest 3.2 with Testing Library; GitHub Actions \\
  \hline
\end{tabular}
\caption{Technology stack, from the dependency manifests}
\label{tab:stack}
\end{table}

\subsection{System Requirements}
\label{sec:sysreq}

FoodLink needs Python 3.10 or later, Node.js with npm (CI uses Node.js~20) and a
web browser; no database server is needed, because migrations run at start-up. A
signing key of at least 32 characters must be set in \code{FOODLINK_SECRET_KEY},
or the backend, CLI and Alembic refuse to start (D-22); no hardware
specification has been defined or measured.

\subsection{Functional Requirements}
\label{sec:fr}

Table~\ref{tab:fr} lists the implemented functional requirements. \emph{[API]}
marks behaviour that is implemented and tested in the API but not yet offered in
the portal.

\begin{table}[!htbp]
\centering
\small
\begin{tabular}{|c|L{12.1cm}|}
  \hline
  \textbf{ID} & \textbf{Requirement (status)} \\
  \hline
  FR-01 & Self-registration as donor, organisation or courier; sign-in; own profile and password \\
  \hline
  FR-02 & A donor posts food details, an optional photo, a pickup pin and a future deadline \\
  \hline
  FR-03 & Posting ranks eligible organisations and records \code{MATCHED} if any exist \\
  \hline
  FR-04 & Verified organisations browse open donations with a live match explanation; readers can fetch the ranked list \emph{[list: API, mobile]} \\
  \hline
  FR-05 & A verified organisation accepts an open donation; overdue donations are refused \\
  \hline
  FR-06 & A courier sets availability, claims a pickup (single winner) and records pickup and delivery \\
  \hline
  FR-07 & The accepting organisation confirms receipt \\
  \hline
  FR-08 & A donor cancels before pickup; the accepting organisation releases a claimed pickup \emph{[API]} \\
  \hline
  FR-09 & Organisations manage standing requirements; donors browse the needs of verified organisations \\
  \hline
  FR-10 & Organisations edit their profile; a name or map-pin change voids verification \emph{[pin: API]} \\
  \hline
  FR-11 & Administrators verify or revoke organisations, run the expiry sweep, export donations as CSV, and manage accounts \emph{[accounts: API]} \\
  \hline
  FR-12 & Every transition is recorded with actor and server time; metrics derive from these records \\
  \hline
\end{tabular}
\caption{Functional requirements}
\label{tab:fr}
\end{table}

\subsection{Non-Functional Requirements}
\label{sec:nfr}

\textbf{Security}, \textbf{privacy} and \textbf{integrity} dominate: every
request is authenticated and checked server-side for role, ownership, state and
trust; reads are scoped in SQL and pickup locations and distances are withheld
from readers who do not need them; and status changes are conditional writes
recorded in an append-only, UTC-stamped ledger. The system must also be
\textbf{explainable} (published weights, per-criterion reasons),
\textbf{abuse-resistant} (rate limits, bounded inputs), \textbf{usable} (short
forms, photo resizing, readable errors) and \textbf{maintainable} (typed
schemas, rules expressed as data, a decision log and 651 automated tests).

\subsection{Database Design}
\label{sec:database}

Six tables, created by one Alembic revision (\code{ae4636b1e6d4}), carry the
workflow (Table~\ref{tab:entities}). Attributes and multiplicities appear in the
class diagram (Figure~\ref{fig:class}).

\begin{table}[!htbp]
\centering
\small
\begin{tabular}{|L{2.6cm}|L{10.7cm}|}
  \hline
  \textbf{Table} & \textbf{Role in the workflow} \\
  \hline
  \code{users} & Every account, with \code{role} as an enumeration; at most one recipient profile and one courier profile each \\
  \hline
  \code{recipients} & Organisation profile: optional coordinates, capacity in meals, \code{is_verified}, acceptance and completion counters \\
  \hline
  \code{volunteers} & Courier profile: base location, availability, completed deliveries \\
  \hline
  \code{donations} & Food details, exact pickup pin, absolute deadline, \code{status}, frozen \code{match_score}; recipient and courier are null until accepted and claimed \\
  \hline
  \code{status_events} & Append-only history: previous and new state, actor, note and server time \\
  \hline
  \code{requirements} & An organisation's standing needs; \code{is_active} is their whole lifecycle \\
  \hline
\end{tabular}
\caption{Database tables and their role}
\label{tab:entities}
\end{table}

Events replace timestamp columns because they record who acted and allow a
state to be entered more than once (D-01). Reliability is computed, not stored:
85 until three acceptances, then completed divided by accepted (D-15); the frozen
\code{match_score} is the deliberate stored exception (D-30). Counters change in
the same transaction as their transition. Invariants live in application code,
since SQLite does not enforce foreign keys by default.

\subsection{API and Backend Design}
\label{sec:api}

The API exposes 27 JSON routes under \code{/api}, with camelCase fields and
OpenAPI documentation at \code{/docs}; Table~\ref{tab:api} summarises the
resource groups.

\begin{table}[!htbp]
\centering
\small
\begin{tabular}{|L{2.3cm}|L{6.8cm}|L{4.2cm}|}
  \hline
  \textbf{Group} & \textbf{Representative routes} & \textbf{Callers} \\
  \hline
  Auth & \code{POST /auth/register}, \code{POST /auth/login}, \code{/auth/me}, \code{POST /auth/password} & Anonymous (rate-limited); any signed-in user \\
  \hline
  Donations & \code{POST /donations}, \code{GET /donations}, \code{/donations/{id}}, \code{/donations/{id}/matches} and \code{/status} & Donor or admin to post; scoped readers; roles per target state \\
  \hline
  Organisations & \code{/recipients}, \code{/recipients/me}, \code{/requirements}, \code{/requirements/{id}}, \code{/volunteers}, \code{/volunteers/me} & Scoped by role and ownership \\
  \hline
  Metrics, Admin & \code{GET /metrics}; \code{/admin/users}, \code{/admin/recipients/{id}/verify}, \code{/admin/maintenance/expire} & Any signed-in user; administrators only \\
  \hline
\end{tabular}
\caption{API resource groups (paths relative to \texttt{/api}; plus \texttt{/api/health})}
\label{tab:api}
\end{table}

Every error carries a human-readable sentence that the frontend shows verbatim
(D-18). Status codes are consistent: 401 means an invalid token, 403 an
insufficient role or trust, 404 not found \emph{or not visible}, 409 a state
conflict, 422 invalid input and 429 a rate limit, and a denied read never
confirms that a record exists (D-24, D-26). In a PATCH an omitted field is
untouched and \code{null} clears only nullable columns (D-61), and role and
verification are absent from self-service schemas (D-16).

\subsection{Frontend Architecture}
\label{sec:frontend}

\code{lib/api.ts} is the only network module; it attaches the token, surfaces
the server's error sentence and signs the user out on a 401 (D-12).
\code{AuthContext} keeps the token in \code{localStorage}, and \code{AppContext}
loads up to 500 donations plus role-specific data, reloading after every write.
Four portals (\code{/donor}, \code{/ngo}, \code{/volunteer}, \code{/admin}) sit
behind route guards that are a convenience only, because the API enforces access
(D-14), and phone-layout screens under \code{/m/*} exist but nothing links to
them (D-20).

\section{Implementation Details}
\label{sec:implementation}

\subsection{Module 1: Core Functionality --- Donation Lifecycle}
\label{sec:module-lifecycle}

A donation passes through nine states: \code{AVAILABLE}, \code{MATCHED},
\code{ACCEPTED}, \code{VOLUNTEER_ASSIGNED}, \code{PICKED_UP}, \code{DELIVERED},
\code{COMPLETED}, \code{CANCELLED} and \code{EXPIRED}. The rules are data
(D-02): \code{ALLOWED_TRANSITIONS} lists the legal moves, and anything else
returns 409, while \code{TRANSITION_ROLES} lists who may cause each target state,
and anyone else receives 403. Pickup, delivery, completion, cancellation and
release also re-read the donation through the caller's own read scope, so a
non-party receives 404 (D-34, D-35); Table~\ref{tab:transitions} summarises them.

\begin{table}[!htbp]
\centering
\small
\begin{tabular}{|L{4.3cm}|L{3.3cm}|L{5.7cm}|}
  \hline
  \textbf{Transition} & \textbf{Who} & \textbf{Key rules and side effects} \\
  \hline
  (new) $\to$ \code{AVAILABLE} & Donor, admin & Future deadline; donor rate limits; bounded input \\
  \hline
  $\to$ \code{MATCHED} & Automatic at posting & At least one eligible organisation; top score frozen \\
  \hline
  \code{AVAILABLE}/\allowbreak\code{MATCHED} $\to$ \code{ACCEPTED} & Verified organisation (for itself), admin & Refused if overdue (409); counts an acceptance; freezes its score \\
  \hline
  \code{ACCEPTED} $\to$ \code{VOLUNTEER_ASSIGNED} & Courier & Conditional claim; the loser receives 409 \\
  \hline
  \code{VOLUNTEER_ASSIGNED} $\to$ \code{ACCEPTED} & Accepting organisation, admin & Release: courier cleared; not a new acceptance (D-41) \\
  \hline
  $\to$ \code{PICKED_UP} $\to$ \code{DELIVERED} & Assigned courier, admin & Owned by the courier bound to the run \\
  \hline
  \code{DELIVERED} $\to$ \code{COMPLETED} & Accepting organisation, admin & Increments completion counters \\
  \hline
  $\to$ \code{CANCELLED} & Owning donor before pickup; admin also after & Withdraws a counted acceptance (D-58) \\
  \hline
  $\to$ \code{EXPIRED} & Admin & Status route, or the sweep over overdue \code{AVAILABLE}/\allowbreak\code{MATCHED} donations \\
  \hline
\end{tabular}
\caption{Lifecycle transitions and enforced rules}
\label{tab:transitions}
\end{table}

Donors may post 10 donations per hour per account and 30 per client address,
with administrators exempt (D-59); free text is capped at 2{,}000 characters, and
\code{imageUrl} must be an image data URL of at most 262{,}144 characters or an
http(s) link (D-56, D-60).

A read-check-write sequence once let two organisations both accept the same
donation; this lost update was reproduced on SQLite before being
fixed~\cite{kleppmann2017}. Every status change is now one statement,
\code{UPDATE donations SET status = :to WHERE id = :id AND status = :from}. A row
count other than one returns 409 and rolls back all side effects (D-55), and the
courier claim binds the courier the same way (D-28).
\code{SELECT ... FOR UPDATE} was rejected because SQLite ignores it.

The expiry sweep is triggered manually and closes overdue \code{AVAILABLE} and
\code{MATCHED} donations with actor-less events. Each \code{StatusEvent} records
the previous and new state, actor, optional note and server time, and drives both
user timelines and platform metrics.

\subsection{Module 2: Core Functionality --- Recipient Matching and Explainability}
\label{sec:module-matching}

\code{matching.py} ranks organisations with a transparent weighted-sum
heuristic; it is \emph{not} a learned model (D-05). An organisation is
excluded entirely, rather than scored low, if it is
unverified, has no coordinates, or lies beyond the radius $R$ (8\,km by default)
(D-06). Let $d$ be the haversine distance in km~\cite{sinnott1984haversine}, $q$
the quantity, $c$ the capacity in meals, $r=q/c$, $m$ the minutes to the
deadline, and $a$ and $k$ the organisation's accepted and completed counts. Each
criterion in Table~\ref{tab:criteria} scores 0--100, and
\begin{equation}
  S = \operatorname{round}\left(0.25\,S_D + 0.25\,S_Q + 0.20\,S_C + 0.15\,S_T + 0.15\,S_R\right).
  \label{eq:overall}
\end{equation}

\begin{table}[!htbp]
\centering
\small
\begin{tabular}{|L{3.3cm}|c|L{8.3cm}|}
  \hline
  \textbf{Criterion} & \textbf{Weight} & \textbf{Score} \\
  \hline
  Distance $S_D$ & 0.25 & $100(1-d/R)$; 0 at or beyond $R$ \\
  \hline
  Quantity fit $S_Q$ & 0.25 & Meals only: $40+60r$ if $r\le1$, else $\max(0,\,100-120(r-1))$ \\
  \hline
  Capacity headroom $S_C$ & 0.20 & Meals only: $100\min(1,(c-q)/100)$; 0 if no room remains \\
  \hline
  Deadline comfort $S_T$ & 0.15 & $100\min(1,(m-t)/120)$ with travel time $t=60d/20$ minutes \\
  \hline
  Reliability $S_R$ & 0.15 & 85 below three acceptances, else $100k/a$ \\
  \hline
\end{tabular}
\caption{Matching criteria and published weights}
\label{tab:criteria}
\end{table}

For units other than meals, $S_Q=S_C=50$ for every candidate, so size cannot
reorder the ranking, and the explanation says size was not assessed (D-42).
Candidates are ordered by $S$, then by how well one of their active standing
requirements fits the donation (same unit only), then by urgency, and finally by
id (D-52). Requirements break ties but never change a score, which is visible to
readers who may not read requirements. Every result carries
the five sub-scores and reasons from fixed thresholds, such as ``2.0 km away in a
straight line --- well inside the collection radius''.

As a worked example, for 60 meals due in three hours the repository's
\code{score_pair} scores a new kitchen 2\,km away with capacity 100 as
$0.25(75)+0.25(76)+0.20(40)+0.15(100)+0.15(85)=73.5$, rounded to 74; a kitchen
5\,km away with capacity 250 and 9 of 10 completions scores 72, and unverified or
9\,km-distant kitchens are excluded.

The frozen \code{matchScore} records the decision: it is written at posting and
replaced at first acceptance, whereas \code{viewerMatch} is an organisation's own
live score for an open donation (D-30). For readers other than that organisation
or an administrator, positions are snapped to a 0.01$^\circ$ grid and
\code{distanceKm} and the frozen score are withheld, which prevents triangulating
private addresses (D-45, D-47); eligibility is still decided on the true
position.

\subsection{Module 3: Core Functionality --- Authentication, Authorization and Trust}
\label{sec:module-auth}

Self-registration refuses the administrator role at the schema level (D-04), and
an organisation account starts \emph{unverified}. Passwords of 8--128 characters
are stored as bcrypt hashes~\cite{provos1999bcrypt}, and
sign-in returns an HS256 JSON Web Token valid for 720
minutes~\cite{rfc7519}. Every request re-reads the user row and ignores the
token's role claim, so a suspension takes effect immediately (D-03). Lockout
guards stop administrators from demoting themselves or the last active
administrator.

Every operation passes four authorization layers: \textbf{role} (403);
\textbf{ownership}, as a SQL clause (404 or an empty list); \textbf{lifecycle
legality} (409); and \textbf{trust}, meaning administrator verification. FoodLink
also controls \emph{what} a record reveals to each reader
(Table~\ref{tab:scopes}).

\begin{table}[!htbp]
\centering
\small
\begin{tabular}{|L{3.4cm}|L{1.4cm}|L{2.2cm}|L{3.3cm}|L{2.4cm}|}
  \hline
  \textbf{Data} & \textbf{Admin} & \textbf{Donor} & \textbf{Organisation} & \textbf{Courier} \\
  \hline
  Donations readable & All & Own & Own, plus the open pool if verified & Unclaimed pickups and own runs \\
  \hline
  Pickup pin, address, donor & All & Own & All it can read & Own claimed runs only \\
  \hline
  Frozen score, true distance & All & None & Own organisation & None \\
  \hline
  Contact directories (phones) & All & None & Own row; own couriers & None \\
  \hline
  Standing requirements & All & Verified organisations' needs & Own & None \\
  \hline
\end{tabular}
\caption{Reader-specific read and disclosure scopes}
\label{tab:scopes}
\end{table}

Verification is the trust boundary: unverified organisations are not ranked,
cannot see the open pool and cannot accept (D-37, D-53). It is read on every
request, so revoking it takes effect at once, and a real change to an
organisation's name, latitude or longitude clears it in the same commit (D-54).
Couriers see only a coarse \code{pickupArea} until a claim binds the pickup to
them (D-57).

\subsection{Module 4: Secondary Features --- Standing Requirements}
\label{sec:module-requirements}

An organisation posts needs with a food type, quantity and unit, beneficiaries,
urgency, a recurrence flag and notes. One flag, \code{is_active}, is their whole
lifecycle (D-29): false retires a need but keeps the row, true reopens it.
Ownership is enforced in the query, so another
organisation's need returns 404, and while donors see active needs of verified
organisations on the \emph{Needs Board}, couriers see none (D-44, D-46). In
ranking, requirements only break ties and add a reason.

\subsection{Module 5: Secondary Features --- Administration and Platform Metrics}
\label{sec:module-admin}

All administration routes sit behind one role dependency. Administrators verify
or revoke organisations, filter and export donations as CSV, run the expiry
sweep, and view the courier roster and analytics; account management is API-only.
They also appear in every entry of \code{TRANSITION_ROLES}, so they can act on a
party's behalf under the same rules. \code{GET /api/metrics} derives volume
totals and, from the event ledger, median \textbf{time-to-claim} (creation to
first \code{ACCEPTED}), median \textbf{handover time} (\code{ACCEPTED} to
\code{DELIVERED}), the \textbf{rescue rate} (completions by the deadline over
completed plus expired donations) and the \textbf{expiry-loss rate} (expired over
that same total). Role impact pages use the account's own donations (D-32).

\subsection{Module 6: Secondary Features --- Role Portals and User Workflows}
\label{sec:module-portals}

Donors post on \emph{Create Food Donation} (geolocation or typed coordinates, as
there is no map picker), follow \emph{My Donations} and browse the \emph{Needs
Board}. Verified organisations use \emph{Available Donations}, \emph{Accepted
Deliveries} and \emph{Food Requirements \& Demand Signals}; couriers use
\emph{Pickup Tasks}, and administrators the \emph{Partner Organizations
Directory} and \emph{Platform Donation Ledger}.

\subsection{Security and Privacy Considerations}
\label{sec:security}

The controls above address this domain's common web-application
risks~\cite{owaspTop10} --- broken access control, privilege escalation,
credential attacks, location exposure, tampered history and abusive posting ---
but residual risks remain: no token revocation, refresh token, MFA or content
security policy; registration reveals whether an email exists and transition
errors reveal a donation's status; photos, descriptions and notes are bounded but
not scoped; and the radius gate is a rate-limited membership oracle.
Section~
ef{sec:limitations} lists the remaining gaps.
